% =================================================
% 2. CONFIGURATION GENERALE
% =================================================

\setstretch{1.15}

% Palette : toutes les couleurs du carnet viennent d'ici
\definecolor{accent}{RGB}{31,78,121}     % titres, liens, définitions
\definecolor{alerte}{RGB}{166,48,40}     % encadrés « À retenir »
\definecolor{methode}{RGB}{36,110,80}    % encadrés « Méthode »
\definecolor{discret}{RGB}{110,110,110}  % code, exemples, éléments secondaires
\definecolor{fondcode}{RGB}{246,246,246}

\hypersetup{
    colorlinks=true,
    linkcolor=accent,
    citecolor=methode,
    urlcolor=accent,
    pdftitle={MiniNumpy - Carnet de mathématiques},
    pdfauthor={Louis Subtil}
}

% Table des matières : chapitres et sections seulement
\setcounter{tocdepth}{1}

% Nombres et unités à la française : \num{0.5} -> 0,5 ; \qty{21}{\ms} -> 21 ms
\sisetup{locale=FR}
\DeclareSIUnit{\ms}{\milli\second}

% Titres
\titleformat{\chapter}[display]
    {\normalfont\bfseries\color{accent}}
    {\large\MakeUppercase{\chaptertitlename}~\thechapter}
    {4pt}
    {\Huge}
    [\vspace{2pt}{\color{accent}\titlerule[1pt]}]
\titlespacing*{\chapter}{0pt}{-10pt}{24pt}
\titleformat*{\section}{\normalfont\Large\bfseries\color{accent}}
\titleformat*{\subsection}{\normalfont\large\bfseries}
